\documentclass[letterpaper,12pt]{article}
\usepackage[T1]{fontenc}
\usepackage[utf8]{inputenc}
\usepackage[spanish]{babel}
\usepackage{framed}
\usepackage{amsfonts,setspace}
\usepackage{amsmath}
\usepackage{amssymb, amsmath, amsthm}
\usepackage{comment}
\usepackage[most]{tcolorbox}
\usepackage{amssymb}
\usepackage{dsfont}
\usepackage{anysize}
\usepackage{pdfpages}
\usepackage{multicol}
\usepackage{enumerate}
\usepackage{graphicx}
\usepackage{float}
\usepackage[dvipsnames]{xcolor}
\usepackage[left=1.5cm,top=0.9cm,right=1.5cm, bottom=1.4cm]{geometry}
\usepackage{fancyhdr}
\pagestyle{fancy}
\definecolor{blue(pigment)}{rgb}{0.2, 0.2, 0.6}
% Page
\oddsidemargin -0.4in
\textwidth 7in
\topmargin -0.6in
\textheight 9.1in
\parindent 0em
\parskip 1ex

%Teoremas, Lemas, etc.
\theoremstyle{plain}
\newtheorem{teo}{Teorema}
\newtheorem{lem}{Lemma}
\newtheorem{prop}{Proposici\'on}
\newtheorem{cor}{Corolario}
\theoremstyle{definition}
\newtheorem{defi}{Definici\'on}
\newtheorem{eje}{Ejemplo}
\newtheorem{obs}{Observaci\'on}
% fin macros

% Comandos más usados
% Conjuntos
\newcommand{\N}{\mathbb{N}}
\newcommand{\Z}{\mathbb{Z}}
\newcommand{\Q}{\mathbb{Q}}
\newcommand{\R}{\mathbb{R}}
\newcommand{\C}{\mathbb{C}}
\newcommand{\RR}{\mathbb{R} \times \mathbb{R}}
% Flechas
\newcommand{\ssi}{\Longleftrightarrow}
\newcommand{\imp}{\Longrightarrow}
\newcommand{\pmi}{\Longleftarrow}
% Trigonométricas
\newcommand{\sen}{\text{sen}}

\usepackage{versions}
\includeversion{solution}

\let\epsilon\varepsilon

\begin{document}
\def\Pauta{1}
%%%%%ESCRIBIR 1 si quieren ver la Prueba con Pauta%%%%
%%%%%ESCRIBIR 0 si quieren ver sólo la Ayudantía%%%%
	
	\fancyhead[L]{\itshape{Escuela de Ingeniería}}
	\fancyhead[R]{\itshape{Universidad de O'Higgins}}
	
	\begin{minipage}{11.5 cm}
         \vspace{-6mm}
		\begin{flushleft}
			\vspace{1mm}
			\textbf{ING1002-2 Cálculo Diferencial e Integral}\\
			\textbf{Profesor:}  Gonzalo Flores G.  \\
			\textbf{Ayudante:}  Juan I. Riquelme \& Cristian Acevedo R. \\
		\end{flushleft}
	\end{minipage}
	
	\begin{picture}(2,3)
		\put(430,9.5){\includegraphics[scale=0.22]{Imagenes/UOH.png}}
	\end{picture}
	\vspace{-6mm}
	\begin{center}
		\LARGE \bf{Ayudantía 12: Técnicas de Integración} %Teorema Fundamental del Cálculo
	\end{center}
    \vspace{-8.8mm}
    \begin{center}
		9 de octubre de 2025
	\end{center}

\begin{enumerate}[\textbf{P\arabic{enumi}.}]
\item Calcule las siguientes integrales usando descomposición en fracciones parciales
\begin{enumerate}
\item $\displaystyle \int_{4}^{10} \frac{1}{x^2 - 9}dx$

\color{blue}
\textbf{Sol:} \begin{align*}
\int_{4}^{10} \frac{1}{x^2 - 9} \, dx 
&= \int_{4}^{10} \frac{1}{(x - 3)(x + 3)} \, dx \\
&= \int_{4}^{10} \left( \frac{A}{x - 3} + \frac{B}{x + 3} \right) dx \\
&\text{Con: } 1 = A(x + 3) + B(x - 3) \\
&\Rightarrow 1 = (A + B)x + (3A - 3B) \\
&\Rightarrow 
\begin{cases}
A + B = 0 \\
3A - 3B = 1
\end{cases}
\Rightarrow A = \frac{1}{6}, \quad B = -\frac{1}{6} \\
\\
\int_{4}^{10} \frac{1}{x^2 - 9} \, dx 
&= \int_{4}^{10} \left( \frac{1}{6(x - 3)} - \frac{1}{6(x + 3)} \right) dx \\
&= \frac{1}{6} \int_{4}^{10} \left( \frac{1}{x - 3} - \frac{1}{x + 3} \right) dx \\
&= \frac{1}{6} \left[ \ln(x - 3) - \ln(x + 3) \right]_{4}^{10} \\
&= \frac{1}{6} \left( \ln(10 - 3) - \ln(10 + 3) - \ln(4 - 3) + \ln(4 + 3) \right) \\
&= \frac{1}{6} \left( \ln 7 - \ln 13 - \ln 1 + \ln 7 \right) \\
&= \frac{1}{3} \ln 7 - \ln 13
\end{align*}
\color{black}
\item $\displaystyle \int \frac{x - 1}{x^2 + 6x + 8} dx$ 

\textcolor{blue}{
\textbf{Sol:} La función a integrar se puede escribir como
\[
\frac{x - 1}{x^2 + 6x + 8} = \frac{x - 1}{(x + 2)(x + 4)} = \frac{A}{x + 2} + \frac{B}{x + 4}
\]
para ciertos valores \( A \) y \( B \). Multiplicando la expresión anterior por \( (x + 2)(x + 4) \) obtenemos
\[
x - 1 = A(x + 4) + B(x + 2) = (A + B)x + 4A + 2B
\]
de donde igualando coeficientes deducimos que \( 4A + 2B = -1 \) y \( A + B = 1 \). La solución de este sistema es \( A = -3/2 \) y \( B = 5/2 \). Así la integral buscada es
\[
\int \frac{x - 1}{x^2 + 6x + 8}dx = \int \frac{-3/2}{x + 2}dx + \int \frac{5/2}{x + 4}dx
\]
\[
= -\frac{3}{2} \ln(x + 2) + \frac{5}{2} \ln(x + 4) + C.
\]
}

\item $\displaystyle \int \frac{1 + x + x^2}{x^3 - 8} \, dx$

\color{blue}
\textbf{Sol:} primero factorizando el denominador:
\[
x^3 - 8 = (x - 2)(x^2 + 2x + 4).
\]
Descomponemos en fracciones parciales:
\[
\frac{1 + x + x^2}{(x - 2)(x^2 + 2x + 4)} = \frac{A}{x - 2} + \frac{Bx + C}{x^2 + 2x + 4}.
\]
\[
1 + x + x^2 = A(x^2 + 2x + 4) + (Bx + C)(x - 2).
\]
\[
1 + x + x^2 = A x^2 + 2A x + 4A + B x^2 - 2 B x + C x - 2 C.
\]
\[
1 + x + x^2 = (A + B) x^2 + (2A - 2B + C) x + (4A - 2C).
\]
Igualamos coeficientes:
\[
\begin{cases}
A + B = 1 \\
2A - 2B + C = 1 \\
4A - 2C = 1
\end{cases}
\]
resolviendo, obtenemos:
\[
A=\frac{7}{12} \quad
C  = \frac{2}{3}
\quad
B = \frac{5}{12}.
\]
luego la descomposición es:
\[
\frac{1 + x + x^2}{x^3 - 8} = \frac{7/12}{x - 2} + \frac{(5/12) x + 2/3}{x^2 + 2x + 4}.
\]
\[
\int \frac{1 + x + x^2}{x^3 - 8} \, dx = \int \frac{7/12}{x - 2} \, dx + \int \frac{(5/12) x + 2/3}{x^2 + 2x + 4} \, dx.
\]
\[
= \frac{7}{12} \ln(x - 2) + \int \frac{(5/12) x + 2/3}{x^2 + 2x + 4} \, dx + C.
\]
Para la segunda integral, completamos el cuadrado en el denominador:
\[
x^2 + 2x + 4 = (x + 1)^2 + 3.
\]
Sea \( u = x + 1 \), entonces \( du = dx \) y \( x = u - 1 \), luego
\[
\int \frac{\frac{5}{12} u + \frac{1}{4}}{u^2 + 3} \, du = \frac{5}{12} \int \frac{u}{u^2 + 3} du + \frac{1}{4} \int \frac{1}{u^2 + 3} du.
\]
\[
\int \frac{(5/12) x + 2/3}{x^2 + 2x + 4} dx = \frac{5}{24} \ln(x^2 + 2x + 4) + \frac{1}{4\sqrt{3}} \arctan \left( \frac{x + 1}{\sqrt{3}} \right) + C.
\]
finalmente
$$
\int \frac{1 + x + x^2}{x^3 - 8} dx = \frac{7}{12} \ln(x - 2) + \frac{5}{24} \ln(x^2 + 2x + 4) + \frac{1}{4\sqrt{3}} \arctan \left( \frac{x + 1}{\sqrt{3}} \right) + C.
$$
\color{black}
\end{enumerate}

\item Calcule las siguientes integrales usando sustitución trigométrica
\begin{enumerate}
\item $\displaystyle \int \frac{dx}{x\sqrt{1+x^{2}}}$ 

\color{blue}
\textbf{Sol:} Sustitución trigonométrica:
\[
x = \tan \theta, \quad dx = \sec^2 \theta \, d\theta,
\implies 
\sqrt{1 + x^2} = \sqrt{1 + \tan^2 \theta} = \sec \theta.
\]
Reemplazamos
\[
I = \int \frac{\sec^2 \theta \, d\theta}{\tan \theta \cdot \sec \theta} = \int \frac{\sec^2 \theta}{\tan \theta \sec \theta} d\theta = \int \frac{\sec \theta}{\tan \theta} d\theta.
\]
\[
I = \int \csc \theta \, d\theta
=
\int \csc \theta \, d\theta = \ln \left( \tan \frac{\theta}{2} \right) + C.
\]
Luego como \( x = \tan \theta \), entonces $\theta = \arctan x.$ Por lo tanto,
\[
\int \frac{dx}{x \sqrt{1 + x^2}} = \ln \left( \tan \left( \frac{\arctan x}{2} \right) \right) + C.
\]
\color{black}
\item $\displaystyle \int \frac{dx}{\sqrt{x^2 + 2x + 2}}$

\color{blue}
\textbf{Sol:} Completamos el cuadrado del denominador:
\[
x^2 + 2x + 2 = (x + 1)^2 + 1.
\]
\[
u = x + 1, \quad du = dx.
\]
La integral queda:
\[
I = \int \frac{du}{\sqrt{u^2 + 1}}.
\]
Usamos la sustitución trigonométrica
\[
u = \tan \theta, \quad du = \sec^2 \theta \, d\theta,
\]
y además
\[
\sqrt{u^2 + 1} = \sqrt{\tan^2 \theta + 1} = \sec \theta.
\]
\[
I = \int \frac{\sec^2 \theta \, d\theta}{\sec \theta} = \int \sec \theta \, d\theta.
\]
\[
\int \sec \theta \, d\theta = \ln \left( \sec \theta + \tan \theta \right) + C.\]
\[
\sec \theta = \sqrt{1 + \tan^2 \theta} = \sqrt{1 + u^2}, \quad \tan \theta = u.
\]
\[
I = \ln \left( (x+1) + \sqrt{(x+1)^2 + 1} \right) + C.
\]

\color{black}
\item $\displaystyle \int_{0}^{\sqrt{2}} \frac{1}{\sqrt{4 - x^{2}}} dx. $

\color{blue}
\textbf{Sol:} Sea la sustitución
\[
x = 2 \sen \theta \implies dx = 2 \cos \theta \, d\theta.
\]
\[
\sqrt{4 - x^2} = \sqrt{4 - 4 \sen^2 \theta} = \sqrt{4(1 - \sen^2 \theta)} = 2 \cos \theta.
\]
Cambiamos los límites de integración,
\[
x=0 \Rightarrow \sen \theta = 0 \Rightarrow \theta = 0,
\]
\[
x = \sqrt{2} \Rightarrow \sqrt{2} = 2 \sen \theta \Rightarrow \sen \theta = \frac{\sqrt{2}}{2} = \frac{1}{\sqrt{2}} \Rightarrow \theta = \frac{\pi}{4}.
\]
Sustituimos en la integral:
\[
I = \int_0^{\pi/4} \frac{1}{2 \cos \theta} \cdot 2 \cos \theta \, d\theta = \int_0^{\pi/4} 1 \, d\theta = \left. \theta \right|_0^{\pi/4} = \frac{\pi}{4}.
\]
\color{black}

\end{enumerate}

\item Muchos cambios de variables de la forma \( x = f(u) \) donde \( f \) es una función trigonométrica están inspirados por las identidades trigonométricas:
\[
\sen^2(x) + \cos^2(x) = 1, \qquad \tan^2(x) + 1 = \sec^2(x)
\]
Pensando en lo anterior resuelva las siguientes integrales:
\begin{enumerate}
\item $\displaystyle \int \tan^5(x) dx$

\color{blue}
\textbf{Sol:} \begin{align*}
\int \tan^{3}(x)\tan^{2}(x)dx &= \int \tan^{3}(x)(\sec^{2}(x)-1)dx \\
&= \int (\tan^{3}(x)\sec^{2}(x) - \tan^{3}(x))dx \\
&= \int \tan^{3}(x)\sec^{2}(x)dx - \int \tan^{3}(x)dx \\
&= \int \tan^{3}(x)\sec^{2}(x)dx - \int \tan(x)(\sec^{2}(x)-1)dx \\
&= \int \tan^{3}(x)\sec^{2}(x)dx - \int \tan(x)\sec^{2}(x)dx + \int \tan(x)dx
\end{align*}
Para las dos primeras integrales, usamos la sustituci\'on $u = \tan(x)$, $du = \sec^{2}(x)dx$:
$$ = \int u^{3}du - \int u du $$
Para la \'ultima integral, usamos $\int \tan(x)dx = \int \frac{\sin x}{\cos x}dx$. Haciendo la sustituci\'on $v = \cos x$, $dv = -\sin x dx$, obtenemos:
$$ \int \frac{-dv}{v} = -\ln|v| $$
Juntando todo:
$$ = \frac{u^{4}}{4} - \frac{u^{2}}{2} - \ln|v| + C $$
Sustituyendo de vuelta $u = \tan x$ y $v = \cos x$:
$$ = \frac{\tan^{4}x}{4} - \frac{\tan^{2}x}{2} - \ln|\cos x| + C. $$
\color{black}
\item $\displaystyle \int \cos^3(x) dx$

\color{blue}
\textbf{Sol:} 
\[
\int \cos^3(x) \, dx = \int \cos(x)(1 - \sin^2(x)) \, dx \quad \quad \textcolor{blue}{\text{(Recordando que}\ \cos^2(x) + \sen^2(x)=1)}
\]
Sea \( u = \sin(x) \), entonces \( du = \cos(x) dx \), y la integral se convierte en:
\[
\int \frac{\cos(x)}{\cos(x)} \cdot (1 - u^2) \, du = u - \frac{u^3}{3} + C = \sen(x) - \frac{\sen^3(x)}{3} + C
\]
\color{black}
\item $\displaystyle \int \sen^4(x) \cos^3(x) \, dx$

\color{blue}
\textbf{Sol:} Separamos el \(\cos(x)\) para hacer un cambio de variable:
\[
= \int \sen^4(x) \cos^2(x) \cos(x) \, dx
\]
Ahora usamos la identidad \(\cos^2(x) = 1 - \sen^2(x)\):
\[
= \int \sen^4(x) (1 - \sen^2(x)) \cos(x) \, dx
\]
Sea el cambio de variable: \( u = \sen(x) \), entonces \( du = \cos(x)\, dx \):
\[
= \int u^4 (1 - u^2) \, du
= \int (u^4 - u^6) \, du = \frac{u^5}{5} - \frac{u^7}{7} + C
\]
\[
= \frac{\sen^5(x)}{5} - \frac{\sen^7(x)}{7} + C
\]

\color{black}
\end{enumerate}
\end{enumerate}
\end{document}